\begin{table}[!ht]
\centering\scriptsize
\setlength{\tabcolsep}{4pt}
\caption{Step-level detection by split, per detector and monitor regime. Detection neither improves nor collapses out of distribution, which is what a taxonomy-defined monitor should show: the rules encode the signal, not the fault schedule.}
\label{tab:appendix_split_det}
\begin{adjustbox}{max width=\textwidth}
\begin{tabular}{@{}llrcc@{}}
\toprule
Arm & Split & Steps & AUC & Firing rate (\%) \\
\midrule
rules / instrumented & train & 10,186 & 0.755 & 20.4 \\
rules / instrumented & dev & 2,672 & 0.758 & 27.1 \\
rules / instrumented & test-id & 2,703 & 0.779 & 27.3 \\
rules / instrumented & test-ood & 8,614 & 0.789 & 21.5 \\
rules / blackbox & test-id & 491 & 0.822 & 13.2 \\
rules / blackbox & test-ood & 603 & 0.759 & 11.6 \\
supervised / instrumented & train & 4,170 & 0.745 & 67.8 \\
supervised / instrumented & dev & 1,426 & 0.774 & 58.7 \\
supervised / instrumented & test-id & 1,324 & 0.762 & 65.1 \\
supervised / instrumented & test-ood & 4,273 & 0.770 & 63.7 \\
supervised / blackbox & test-id & 487 & 0.752 & 61.2 \\
supervised / blackbox & test-ood & 573 & 0.747 & 61.1 \\
\bottomrule
\end{tabular}
\end{adjustbox}
\end{table}
